%-*-latex-*-
%-*-latex-*-
\newcommand\COURSE{ciss350}
\newcommand\ASSESSMENT{q10602}
\newcommand\ASSESSMENTTYPE{Quiz}
\newcommand\POINTS{\textwhite{xxx/xxx}}

\input{myquizpreamble}
\input{yliow}
\input{\COURSE}
%-*-latex-*-

%-*-latex-*-

%-*-latex-*-
\renewcommand\TITLE{\ASSESSMENTTYPE \ \ASSESSMENT}

\renewcommand\EMAIL{}
\renewcommand\AUTHOR{eagraham1@cougars.ccis.edu}

\textwidth=6in
\begin{document}
\topmatterthree

%------------------------------------------------------------------------------
\nextq
Write a function \verb!f!
that sorts \verb!v! (a \verb!std::vector< int >! object)
by the lowest order digit in the following
way.
(The lower order digit of \verb!42! is \verb!2!.)
Create an array \verb!x! of 10 \verb!std::list! objects.
\verb!x[0]! will be a list of integers with lower order digit \verb!0!.
\verb!x[1]! will be a list of integers with lower order digit \verb!1!.
Etc.
If the integer you read from \verb!v! is \verb!42!, you perform
insert tail on \verb!x[2]! with \verb!42!.
If the integer you read from \verb!v! is \verb!705!, you perform
insert tail on \verb!x[5]! with \verb!705!.
For instance if \verb!42, 62, 32! are the only integers in \verb!v!
with lowest order digit of \verb!2!, then \verb!x[2]! will
be the linked list \verb!42, 62, 32! where \verb!42! is the head.
After all values of \verb!v! are placed in \verb!x!, you copy the values from
\verb!x! back into \verb!v! from index 0 to index 9,
and for each \verb!x[i]!, you process the values from head to tail.

(This is essentially a radix sort on the lowest order digit.
If you repeat the above process for the 10s digit, for the 100s digit, etc.,
you'll get the radix sort using linked lists.)

A skeleton code is given. Here's a test case:
\begin{console}[commandchars=\\\{\},fontsize=\small]
\userinput{52 31 75 73 51 -1}
v: [52, 31, 75, 73, 51]
[]
[31, 51]
[52]
[73]
[]
[75]
[]
[]
[]
[]
v: [31, 51, 52, 73, 75]
\end{console}
\\
\ANSWER
\begin{answercode}
#include <iostream>
#include <string>
#include <vector>
#include <list>

std::ostream & operator<<(std::ostream & cout, const std::vector< int > & v)
{
    std::string delim = "";
    cout << "[";
    for (auto x: v)
    {
        cout << delim << x;
        delim = ", ";
    }
    cout << "]";
    return cout;
}

std::ostream & operator<<(std::ostream & cout, const std::list< int > & v)
{
    std::string delim = "";
    cout << "[";
    for (auto x: v)
    {
        cout << delim << x;
        delim = ", ";
    }
    cout << "]";
    return cout;
}

void f(std::vector< int > & v)
{
    std::vector< std::list< int > > lists(10);

    // TODO: Add values in v to lists
    for (int x : v)
    {
        int digit = x % 10;
        lists[digit].push_back(x);
    }

    for (auto list: lists)
    {
        std::cout << list << '\n';
    }

    // TODO: Copy values in lists to v
    int index = 0;
    for (int i = 0; i < 10; ++i)
    {
        for (int x : lists[i])
        {
            v[index++] = x;
        }
    }
    
}

int main()
{
    std::vector< int > v;
    while (1)
    {
        int t;
        std::cin >> t;
        if (t == -1) break;
        v.push_back(t);
    }
    std::cout << "v: " << v << '\n';
    f(v);
    std::cout << "v: " << v << '\n';

    return 0;
}
\end{answercode}

%------------------------------------------------------------------------------
%-*-latex-*-
\newpage
\input{instructions}
\end{document}

